\documentclass[11pt,a4paper]{article}

\usepackage[T1]{fontenc}
\usepackage[utf8]{inputenc}
\usepackage{lmodern}
\usepackage[margin=2.5cm]{geometry}
\usepackage{microtype}
\usepackage{enumitem}
\usepackage{xcolor}
\usepackage{hyperref}

\setlength{\parindent}{0pt}
\setlength{\parskip}{0.6em}
\setlist[itemize]{topsep=0.2em,itemsep=0.25em}

\definecolor{linkblue}{HTML}{1F4E79}
\hypersetup{
    colorlinks=true,
    linkcolor=linkblue,
    urlcolor=linkblue,
    pdftitle={VIP Week 5 Report},
    pdfauthor={Yorgo Azzi}
}

\begin{document}

\begin{center}
    {\Large\bfseries VIP Week 5 Report}\\[0.5em]
    {\large Yorgo Azzi}\\[0.25em]
    October 14, 2026
\end{center}

\section{Overview}

This week, I implemented the new performance approach proposed in the previous
report. The system now has a second mode in which the Markov model generates
three accompaniment voices while the performer controls the soprano directly
with the right hand. I also combined the camera interface and live score in one
application, added soprano timbre and ornament controls, and improved the
response and continuity of the hand-controlled melody.

The purpose of these changes was to give the performer a more direct musical
role. Instead of requesting chord changes and waiting for the model to apply
them, the performer can now shape an audible melodic line over an accompaniment
that remains generated by the model.

\section{Work Completed}

\subsection{New Soprano Performance Mode}

The application now provides two selectable modes. The original Full SATB mode
keeps the four-voice generator and chord gestures. The new Soprano Performance
mode disables direct chord control and divides the musical roles as follows:

\begin{itemize}
    \item the Markov and voice-leading engine generates alto, tenor, and bass;
    \item the model continues to generate the underlying chord progression;
    \item the right hand performs the soprano independently of the current
    chord, while remaining within the selected key and mode;
    \item the left hand controls the accompaniment register, tempo, and
    articulation.
\end{itemize}

The interface includes buttons for switching modes without restarting and a
Pause/Play button. Pausing stops the sounding notes while hand tracking remains
active, which prevents a burst of queued movement when playback resumes.

\subsection{Direct Right-Hand Soprano Control}

Right-hand vertical movement is converted into successive scale steps. Moving
up produces ascending scale notes, and moving down produces descending scale
notes. This differs from the previous register mapping because it creates a
performable melodic sequence rather than shifting a generated range.

The soprano also supports a small set of gestures:

\begin{itemize}
    \item a horizontal shake produces a trill;
    \item a quick index-finger bend and release produces a mordent;
    \item a pinch sustains the current pitch;
    \item a fist releases the soprano and produces a rest.
\end{itemize}

An earlier arpeggio gesture was removed because it was difficult to recognize
reliably and overlapped with other finger movements. Finger count is now used
to change only the soprano timbre: piano, violin, choir, flute, or synth pad.
This leaves the accompaniment sound unchanged.

The stationary-hand behavior was also improved. The soprano previously
repeated sixteenth notes even when the hand was still. It now holds one
continuous note. If the note crosses a measure boundary, the score uses a tie.
The next vertical movement replaces the held note immediately, preserving a
legato melodic flow without reducing responsiveness.

\subsection{Faster Musical Response}

Soprano control uses 16 timing ticks per beat. A detected scale step is applied
on the next fine tick instead of waiting for the next accompaniment onset,
beat, or measure. At 60 BPM, one tick is 62.5 ms; at 140 BPM, it is about
26.8 ms. The old soprano note is turned off and the new pitch is sent directly
to MIDI on that tick.

This separates the performer's melodic timing from the slower Markov rhythm of
the accompaniment. Alto, tenor, and bass can continue holding their notes while
the user changes only the soprano. Trills and mordents retain short durations,
while a stationary soprano has no automatic MIDI note-off and remains active
until movement, a rest gesture, pause, or mode change ends it.

\subsection{Integrated Camera and Live Score}

The camera and score are now displayed side by side in one PyQt application.
The left panel shows the mirrored camera image, MediaPipe landmarks, detected
gestures, mode controls, and performance status. The right panel shows the live
score and keeps completed four-measure sections available for scrolling.

The score is produced from the same symbolic events sent to MIDI. A background
worker converts these events to MusicXML with \texttt{music21}, and Verovio
engraves the result as SVG. Because this work runs separately, notation does
not delay audio playback.

During testing, Qt initially displayed incomplete SVG graphics, leaving stems
and fragments without noteheads or staff lines. Verovio's reusable SVG symbols
were expanded into ordinary paths before display, which corrected the score.
Dense four-measure sections are also allowed to wrap onto multiple systems for
better readability.

\section{Results}

The main results were:

\begin{itemize}
    \item The user can switch between Full SATB and Soprano Performance modes
    in the same application.
    \item The Markov model produces a continuous three-voice accompaniment
    while the right hand performs an independent scale-based soprano.
    \item Vertical movements can change soprano pitch on the next fine tick
    without retriggering alto, tenor, or bass.
    \item A stationary hand produces a continuous held note instead of repeated
    attacks, while later movement remains immediately effective.
    \item Finger count changes only the soprano instrument, and the trill,
    mordent, pinch, and fist gestures have separate musical functions.
    \item The camera, controls, and complete scrollable score now appear in one
    desktop window.
    \item The notation renderer now displays complete musical symbols and wraps
    dense passages more clearly.
    \item All 51 automated tests passed, covering hand gestures, the two
    performance modes, sustained MIDI notes, fast pitch replacement, Markov
    generation, live notation, and the integrated interface.
\end{itemize}

\section{Challenges}

The interaction is more direct than the previous chord-control approach, but
the perceived latency still includes camera capture, MediaPipe processing, the
wait for the next fine musical tick, and the latency of the MIDI synthesizer.
The theoretical tick delay is small, but end-to-end physical measurements are
still needed.

Gesture recognition also depends on camera angle, lighting, hand orientation,
and tracking stability. Rapid finger movement can make the trill or mordent
harder to recognize, while different General MIDI synthesizers may produce
noticeably different instrument sounds. These controls therefore need further
physical testing and threshold tuning.

\section{Next Steps}

\begin{itemize}
    \item Integrate a Google Magenta model into the realtime generation system
    and determine which model can produce responsive accompaniment from a live
    symbolic-MIDI input.
    \item Use the musical output of Soprano Performance mode as input to the
    Magenta model. Magenta would listen to the material produced by the current
    system and generate an additional accompaniment around the performance.
    \item Add an interface control that selects the information sent to
    Magenta. One option would send only the hand-performed soprano, allowing
    Magenta to create an accompaniment from the user's melody. The second
    option would send the soprano together with alto, tenor, and bass, allowing
    Magenta to respond to the complete texture produced by the performer and
    the Markov model.
\end{itemize}

\section{Resources}

\begin{itemize}
    \item Project repository:
    \url{https://github.com/YorgoAzzi/MarkovModel}
    \item MediaPipe Hand Landmarker documentation:
    \url{https://ai.google.dev/edge/mediapipe/solutions/vision/hand_landmarker}
    \item Verovio documentation:
    \url{https://book.verovio.org/}
\end{itemize}

\end{document}
